% 第五章 神经形态计算初步探索
\chapter{神经形态计算初步探索}

非易失性、低功耗以及丰富的动力学行为，使磁性拓扑结构成为神经形态计算硬件的合适候选。本章承接前几章对 Hopfion 稳定性与动力学调控的讨论，从物理机制映射的视角出发，初步评估其在神经形态计算中的应用前景。


\section{Hopfion 动力学特征与神经元功能的类比}

生物神经元的功能可归纳为四个环节：输入信号的接收、跨越阈值后的脉冲激发、向下游的脉冲传播，以及响应之后的不应期抑制。需要指出的是，第四章给出的仿真结果显示，Hopfion 在自旋波驱动下的行为与上述环节之间存在可被直接对照的物理对应关系：

\begin{enumerate}
  \item \textbf{阈值效应}：从幅度扫描（\cref{subsec:amplitude-scaling}）来看，自旋波幅度 $B_0$ 存在明显的驱动门槛——$B_0 \leq \SI{0.1}{T}$ 时 Hopfion 几乎原地不动，$B_0 \geq \SI{0.5}{T}$ 才会触发有效位移。此类阶跃式响应与神经元"全或无"激发特性高度相似。
  \item \textbf{频率选择性}：频率扫描结果（\cref{subsec:freq-sweep}）表明，Hopfion 对特定频率（典型如 $\SI{1100}{GHz}$）的自旋波响应强烈，而在其他频段，尤其是 $\SIrange{400}{600}{GHz}$ 的死区，几乎保持静止。该选择性可类比为神经元对特定突触输入模式的筛选机制。
  \item \textbf{方向编码}：拓扑 Hall 效应（\cref{subsec:hall-effect}）使 srcX 与 srcZ 激励产生走向相反的运动（$+z$ 与 $-z$），而 Hall 角在拓扑保护下保持高度稳定。从仿真结果来看，这为在单个器件中同时承载兴奋性（$+z$）与抑制性（$-z$）两类信号提供了物理基础。
  \item \textbf{双向可控}：双向控制实验（\cref{subsec:bidirectional-control}）表明，切换自旋波频率或改变传播方向即可引导 Hopfion 的可逆位移，呈现出与神经元激发-恢复循环相似的行为特征。
\end{enumerate}


\section{概念器件设计}

依据上述映射关系，本文构想了一种以 Hopfion 为核心的自旋波神经形态器件，其构成如下：

\begin{enumerate}
  \item \textbf{自旋波输入端}：由微波天线阵列扮演"突触"角色，借助自旋波将输入信号导入 Hopfion。天线间频率、幅度和传播方向的不同组合，共同张成多维输入空间。
  \item \textbf{Hopfion 处理单元}：作为"神经元胞体"的 Hopfion 以位置变化编码输出；其阈值效应天然实现非线性激活，频率选择性则承担输入滤波的职能。
  \item \textbf{位置读出}：利用磁阻效应（TMR/GMR）或拓扑 Hall 效应读取 Hopfion 的位置信息，并转换为电信号输出。
  \item \textbf{权重调节}：把自旋波幅度 $B_0$ 视作突触权重，$v \propto B_0^2$ 的标度律保证了权重到输出之间存在确定性映射。
\end{enumerate}

就优势而言，本方案有以下几点：相较于二维 Skyrmion，Hopfion 的三维结构在存储密度和动力学自由度上更具潜力；拓扑保护赋予信号处理较强的鲁棒性；采用自旋波驱动还能规避电流引起的焦耳热。

当然，从构想走向实现仍有若干门槛需要跨越。实验上 Hopfion 目前仅在低温条件下被观测到\cite{zheng2023hopfion}，室温稳定性尚待突破；在竞争交换体系中，其特征尺寸约 $\SI{2}{nm}$，已贴近当前加工工艺的极限；此外，纳米尺度下自旋波的衰减问题也需要做进一步的研究。


Hopfion 从一个基础物理对象转向神经形态计算硬件候选的可能性，在本章得到了较为正面的初步评估。阈值效应、频率选择性、方向编码以及双向可控——这四项由前几章仿真所刻画的动力学特征，恰好能与神经元的激发、滤波、兴奋/抑制与恢复四个环节建立起较为直接的物理对应。以此映射关系为支点，本文进一步构想了一类以 Hopfion 为核心的自旋波神经形态概念器件。然而要把该构想推向实验验证，材料体系的突破、加工工艺的进展以及室温稳定性问题的解决，仍是缺一不可的前置条件。
